\section{Experiments}
\label{sec:experiments}

We investigate three questions: whether naturally shorter correct traces provide better student supervision; whether length-associated SAE features can control teacher generation; and whether the generated traces change student accuracy and reasoning length.
We report the completed exploratory GSM8K study, separating the initial length comparison from the subsequent intervention experiment.
The evidence covers one teacher--student pair and one SAE training seed; the three student training seeds quantify variation in SFT, not in SAE discovery.

\subsection{Experimental Setup}
\label{sec:experimental_setup}

\paragraph{Models and data.}
We use Qwen2.5-7B-Instruct as the teacher and Qwen2.5-1.5B-Instruct as the student.
The initial training cohort contains 881 fixed GSM8K training problems~\citep{cobbe2021training}.
Teacher sampling uses temperature 0.7, top-$p$ 0.95, and a maximum of 512 generated tokens.
Student evaluation uses the fixed 1,269-problem cohort \texttt{test[50:1319]}, greedy decoding, and the same generation cap; the first 50 test problems are reserved for smoke checks.
Accuracy measures correctness of the extracted final answer, and output length counts generated tokens.
We do not use final evaluation accuracy to select the intervention.

\paragraph{Student training.}
All conditions start from the same pretrained student and use completion-only SFT with LoRA rank 4, scaling factor 16, dropout 0.05, and all linear modules targeted.
We train for one epoch over each constructed dataset, using batch size 4, gradient accumulation 1, learning rate $2\times10^{-5}$, warmup ratio 0.03, BF16, and a maximum sequence length of 2,048 tokens without packing.
Student seeds are 17, 42, and 73.
The intervention study compares equal-example and approximately equal-target-token budgets separately.
Token matching repeats complete solutions rather than truncating them; consequently, it does not equalize optimizer steps or per-question weights.

\paragraph{SAE construction.}
The initial teacher pool contains 16 candidates per problem, or 14,096 traces.
We fit shared TopK SAEs to mixed teacher residual activations, including raw candidates regardless of correctness; correct traces are used for the short--long feature comparison.
Question-level splits contain 617 training, 132 development, and 132 test problems.
The exploratory SAE grid covers layers 10, 17, and 23, TopK sparsities 32 and 64, expansion factor 8, and 1,500 training steps, with SAE seed 17.
The primary layer-17, TopK-64 dictionary yields eight short-associated and eight long-associated features under the registered screening procedure.
These labels describe associations with trace length, not verified reasoning operations.

\subsection{Natural Trace Length and Student Learning}
\label{sec:natural_length}

For each problem, we require at least four distinct correct candidates that do not reach the token cap.
We choose short and long representatives from the bottom and top 20\% length bands, respectively, and a medium representative near the median.
All three student conditions retain the same 881 problems and use equal-example training.
Mean accuracy across student seeds is 66.67\% for short, 66.33\% for medium, and 64.20\% for long traces, compared with the independently evaluated base accuracy of 67.45\% in this initial run.
The short--long difference is 2.47 percentage points, but does not pass Holm correction across the three registered comparisons ($p=0.0852$).
Thus, this sample provides a directional short--long difference without establishing either a corrected significant advantage or an improvement over the base student.
Because supervision volume differs, this comparison alone cannot attribute the ordering to redundant content.

\subsection{Feature Intervention and Generation Controls}
\label{sec:teacher_intervention}

We screen short-feature enhancement and joint short-enhancement/long-direction subtraction at relative perturbation magnitudes $\rho\in\{0.05,0.15,0.30\}$.
Each target condition has a random control matched in decoder-column count and actual perturbation norm; target features are excluded from the random sets.
Together with unmodified generation, this gives 13 development conditions.
The joint intervention uses a fixed composite direction, rather than activation-gated suppression of individual long features.
Interventions are applied at the layer-17 post-block residual position used to predict each next token, beginning with the first generated token.
Conditions share per-question random uniforms and prompts but maintain independent attention caches.

We use 64 development problems to select the shortest eligible condition.
Eligibility requires at least 5\% mean shortening, no more than a five-percentage-point decline in empirical accuracy, a paired 95\% bootstrap length-difference interval below zero, and a shorter mean output than the matched random control.
The selected short-enhancement condition has $\rho=0.30$ and uses eight short-associated decoder directions.
On a subsequent 64-problem check, it reduces mean length from 246.44 to 209.23 tokens, with both unmodified and targeted generation answering 64/64 correctly.
These cohorts were previously observed during SAE analysis, so this check is exploratory; its accuracy threshold is not a non-inferiority test.

After fixing the intervention, we generate four candidates for every original problem under unmodified, targeted, and matched-random generation, producing 10,572 outputs.
Table~\ref{tab:teacher_generation} reports the entire candidate pool, including incorrect and capped responses.
Targeted generation reduces mean length by 13.27\% relative to unmodified generation.
A question-level bootstrap, retaining the four candidates together within each problem, gives a mean length difference of $-33.20$ tokens with a 95\% interval of $[-35.24,-31.22]$.
The correctness difference is $+0.20$ percentage points with interval $[-0.43,+0.79]$.
These full-pool statistics are post-hoc descriptive summaries and do not establish either improved correctness or population-level accuracy preservation.

\begin{table}[t]
\centering
\small
\begin{tabular}{lrrr}
\hline
Generation condition & Correct (\%) & Mean tokens & Cap hits \\
\hline
Unmodified & 98.69 & 250.27 & 12 \\
SAE short enhancement & 98.89 & 217.07 & 1 \\
Matched random & 97.70 & 293.02 & 38 \\
\hline
\end{tabular}
\caption{Complete teacher candidate pools: 3,524 outputs per condition on the same 881 problems. Cap hits count responses reaching the 512-token limit.}
\label{tab:teacher_generation}
\end{table}

\subsection{Distillation from Intervention-Generated Traces}
\label{sec:student_results}

\paragraph{Comparisons and common support.}
For each newly generated condition, we deduplicate responses and select the shortest correct, non-capped solution per problem.
The unmodified, SAE, and random conditions cover 881, 880, and 879 problems, respectively.
Intersecting their support with the initial natural-short dataset yields 878 problems, on which all four conditions are retrained.
The natural-short comparator uses the initial 16-candidate quantile-band selection, whereas the three new conditions use shortest-correct selection from four candidates.
It therefore compares data-construction pipelines with different sampling and selection rules, rather than isolating the effect of intervention.

\paragraph{Budgets and statistical analysis.}
Four data conditions, two budgets, and three seeds yield 24 adapters, plus a separately evaluated base student.
Equal-example datasets contain 878 records each and require 220 optimizer steps.
Their target-token totals are 188,540 for unmodified, 165,530 for SAE, 215,643 for random, and 188,420 for natural short.
The equal-token construction brings all totals into the range 215,643--215,722, with 252, 286, 220, and 251 optimizer steps, respectively.
We report seed-mean accuracy and 95\% intervals from 10,000 crossed bootstrap replicates over training seeds and paired evaluation problems.
Holm correction covers six registered SAE-versus-control accuracy comparisons across the two budgets.
The random condition has identical data, weights, and predictions across budgets for each seed; these entries are not independent replications.

\begin{figure}[t]
\centering
\includegraphics[width=\linewidth]{../../results/ncsu_sae_intervention_v1/exploratory/student_followup/analysis/student_comparison.pdf}
\caption{Student accuracy in the completed intervention study. Bars show means over three training seeds; dots show individual seeds, not confidence intervals. The dashed line marks the base student evaluated in this run. Natural-short data use a different candidate pool and selection rule from the three intervention-generation conditions.}
\label{fig:student_accuracy}
\end{figure}

\paragraph{Accuracy.}
Figure~\ref{fig:student_accuracy} summarizes the four conditions.
With equal examples, SAE-generated supervision gives 70.84\% accuracy, compared with 69.92\% for unmodified outputs, 68.87\% for random intervention, and 66.93\% for natural-short supervision.
With equal target tokens, the corresponding accuracies are 70.34\%, 70.42\%, 68.87\%, and 66.96\%.
The base student in this run achieves 66.51\%; the earlier run's 67.45\% base is not substituted for this baseline.
The source of the between-run base difference has not been isolated.

The SAE-minus-unmodified difference is $+0.92$ percentage points with 95\% interval $[-1.34,+3.23]$ under equal examples, and $-0.08$ with interval $[-2.18,+1.97]$ under equal tokens.
Neither comparison establishes an accuracy improvement; the Holm-adjusted $p$-values are 0.8432 and 0.9654.
Comparisons against the matched-random condition are also inconclusive after correction.
SAE outperforms the natural-short comparator under both budgets (adjusted $p=0.0060$ and $0.0110$), but the sampling and selection differences prevent attribution of these gaps solely to SAE steering.

\paragraph{Student output length.}
Under equal examples, mean student output length decreases from 277.14 tokens with unmodified supervision to 245.98 with SAE supervision, a descriptive reduction of 11.24\%.
Under equal tokens, it decreases from 276.55 to 245.31 tokens, or 11.30\%.
The base student averages 239.64 tokens, so the observed shortening is relative to unmodified distillation rather than to the pretrained student.
The results support transfer of shorter-output behavior in this setting, while leaving the accuracy benefit unresolved.
Output-token reductions alone do not measure end-to-end savings, which must also account for SAE fitting, candidate generation, selection, and intervention overhead.

\subsection{Additional Comparisons and Scope}
\label{sec:experiment_scope}

The completed results do not include concise prompting, dense short--long steering, ASC, or difficulty-aware post-hoc compression.
These are important comparisons for determining whether SAE decomposition provides value beyond simpler length control or existing compression methods~\citep{azizi2026activation,wu2025concise}.
A length-matched comparison is also needed before attributing any student advantage to supervision content beyond brevity.
The present study does not establish superiority over these methods.

A separately registered multi-answer extension compares one, two, and four distinct correct responses per problem under SAE and unmodified generation, retaining the same 878 problems.
Repeating the shortest response two or four times provides controls for record count, question exposure, and optimizer steps.
Distinct text does not guarantee a distinct mathematical solution, and these controls do not match target-token counts.
The extension has not produced a complete audited result set at the time of this draft; we report its design without drawing conclusions from partial runs.

Finally, the study uses previously observed feature-analysis cohorts and a single SAE seed, and lexical or answer-format explanations remain possible.
Matched random directions control feature count and perturbation magnitude, not feature semantics or activation prevalence.
The current evidence therefore supports an exploratory account of length control and its downstream effects on GSM8K, rather than a general claim that long traces are redundant or that SAE steering improves student accuracy.
